% Minimal preamble for Open Logic Project sources (pinned commit 1e960beff9ed7835bf3e3f1335e21af3439cd107).
%
% OLP's own master preamble cannot run on the project host: TeX Live 2025 there lacks
% prooftree.sty, ntheorem.sty, extarrows.sty, mathdots.sty, isomath.sty and subfiles.sty.
% Following the project's own rule ("if the author preamble needs packages the toolchain lacks,
% do not invent macros -- use a minimal working preamble"), this file declares, verbatim in
% spirit and numbering behaviour, exactly the theorem-like environments OLP defines in
% sty/open-logic-book-envs.sty (thm/lem/prop/cor/defn/ex/prob/rem/note/case/axiom/conv),
% plus the shorthand accents OLP uses (`!!a`, `!!b`, `!!{...}`).
%
% Author text is untouched; only the packaging is ours.
\documentclass[10pt,a4paper]{article}

\usepackage{amsmath}
\usepackage{amssymb}
\usepackage{amsthm}
\usepackage{graphicx}
\usepackage{hyperref}
\usepackage{enumitem}
\usepackage[framemethod=default]{mdframed}

% ---- OLP shorthand accents (from sty/open-logic.sty) ----
% NB: \! is taken by LaTeX (negative thin space), so we only define the !!a / !!b families
% and the \olup helper used by them.
\newcommand*{\olup}[1]{\MakeUppercase{#1}}
\def\!!!a#1{\olup{#1}}
\def\!!a#1{\olup{#1}}
\def\!!b#1{\olup{#1}}

% ---- theorem-like environments, same family numbering as OLP ----
\theoremstyle{plain}
\newtheorem{thm}{Theorem}[section]
\newtheorem{lem}[thm]{Lemma}
\newtheorem{prop}[thm]{Proposition}
\newtheorem{cor}[thm]{Corollary}
\newtheorem{axiom}{Axiom}

\theoremstyle{definition}
\newtheorem{defn}[thm]{Definition}
\newtheorem{ex}[thm]{Example}
\newtheorem{prob}{Problem}[section]
\newtheorem{probd}{Problem}
\newtheorem{conv}[thm]{Convention}

\theoremstyle{remark}
\newtheorem{rem}[thm]{Remark}
\newtheorem{note}[thm]{Note}
\newtheorem{case}[thm]{Case}

% ---- environments OLP defines as wrappers ----
\newenvironment{probtag}[1]{\begin{prob}}{\end{prob}}
\newenvironment{explain}{}{}
\newenvironment{explain*}{}{}
\newenvironment{editorial}{}{}
\newenvironment{digress}{}{}
\newenvironment{digress*}{}{}
\newenvironment{history}{}{}
\newenvironment{reference}{}{}
\newcommand*{\olfileid}[3]{}
\newcommand*{\olsection}[1]{\section*{#1}}
\newcommand*{\olchapter}[3]{}
\newcommand*{\olpart}[2]{}
\newcommand*{\olimport}[1]{}
\newcommand*{\olinclude}[1]{}
\newcommand*{\olpath}[1]{}
\newcommand*{\ollangid}[1]{}
\newcommand*{\problemtag}[1]{}
\newcommand*{\problemauthor}[1]{}
\newcommand*{\ollabel}[1]{\label{#1}}
% \olref takes 0-3 OPTIONAL arguments plus one mandatory label:
%     \olref{lab}  |  \olref[from]{lab}  |  \olref[from][to]{lab}  |  \olref[a][b][c]{lab}
% A single-argument definition silently swallows "[sth]" and emits a broken reference
% (measured on the host: the text came out as "??iso]x"). The chain below consumes the
% brackets and cross-references only the mandatory argument.
\makeatletter
\newcommand{\olref}{\@ifnextchar[{\olref@i}{\olref@final}}
\def\olref@i[#1]{\@ifnextchar[{\olref@ii}{\olref@final}}
\def\olref@ii[#1]{\@ifnextchar[{\olref@iii}{\olref@final}}
\def\olref@iii[#1]{\olref@final}
\def\olref@final#1{\ref{#1}}
\makeatother
\newcommand*{\olnewpage}{\clearpage}

% ===================== OLP compatibility layer (1) =====================

\makeatletter

% ---------------- proof trees (OLP's prooftree is a bussproofs derivative) ----------------
\providecommand{\AxiomC}[1]{}
\providecommand{\UnaryInfC}[1]{}
\providecommand{\BinaryInfC}[1]{}
\providecommand{\TrinaryInfC}[1]{}
\providecommand{\QuaternaryInfC}[1]{}
\providecommand{\QuinaryInfC}[1]{}
\providecommand{\RightLabel}[1]{}
\providecommand{\LeftLabel}[1]{}
\providecommand{\doubleLine}{}
\providecommand{\noLine}{}
\providecommand{\DisplayProof}{}
\providecommand{\fCenter}{}
\newenvironment{prooftree}{}{}
\newenvironment{ltprooftree}{}{}
\newenvironment{oltableau}{}{}

% ---------------- logical notation ----------------
\providecommand{\lif}{\mathbin{\rightarrow}}
\providecommand{\lforall}{\mathord{\forall}}
\providecommand{\lexists}{\mathord{\exists}}
\providecommand{\ltrue}{\top}
\providecommand{\lfalse}{\bot}
\providecommand{\Ltrue}{\top}
\providecommand{\Lfalse}{\bot}
\providecommand{\Sequent}{\mathrel{\Rightarrow}}
\providecommand{\Proves}{\mathrel{\vdash}}
\providecommand{\Entails}{\mathrel{\models}}
\providecommand{\Subst}[3]{#1[#2/#3]}
\providecommand{\sFmla}[2]{#1(#2)}
\providecommand{\pFmla}[2]{#1(#2)}
\providecommand{\formula}{\relax}
\providecommand{\True}{\top}
\providecommand{\False}{\bot}
\providecommand{\Args}[1]{#1}
\providecommand{\Ts}{\ensuremath{\top}}
\providecommand{\Fi}{\ensuremath{\bot}}
\providecommand{\lfalsum}{\bot}

% ---------------- objects, structures, metatheory ----------------
\providecommand{\Obj}{}
\providecommand{\Struct}{\ensuremath{\mathfrak{A}}}
\providecommand{\Domain}{\ensuremath{|\mathfrak{A}|}}
\providecommand{\Assign}{\ensuremath{s}}
\providecommand{\Sat}[2]{\ensuremath{#1 \models #2}}
\providecommand{\mSat}[3]{\ensuremath{#1 \models #2[#3]}}
\providecommand{\mModel}[2]{\ensuremath{#1 \models #2}}
\providecommand{\Lang}{\ensuremath{\mathcal{L}}}
\providecommand{\Th}{\ensuremath{\mathrm{Th}}}
\providecommand{\Log}{\ensuremath{\mathrm{L}}}
\providecommand{\tuple}[1]{\langle #1 \rangle}
\providecommand{\num}[1]{\underline{#1}}
\providecommand{\cfind}[1]{\varphi_{#1}}
\providecommand{\Nat}{\ensuremath{\mathbb{N}}}
\providecommand{\Mult}{\ensuremath{\mathrm{Mult}}}
\providecommand{\Add}{\ensuremath{\mathrm{Add}}}
\providecommand{\Exp}{\ensuremath{\mathrm{Exp}}}
\providecommand{\Zero}{\ensuremath{0}}
\providecommand{\Succ}{\ensuremath{\mathrm{S}}}
\providecommand{\lambd}{\ensuremath{\lambda}}
\providecommand{\bredpar}{\mathrel{\rightarrow_\beta}}
\providecommand{\bcd}[1]{\ensuremath{\ulcorner #1 \urcorner}}
\providecommand{\Set}[1]{\ensuremath{\{#1\}}}
\providecommand{\indcase}[1]{\par\noindent\emph{Case #1}.\quad}
\providecommand{\indfrm}[1]{\par\noindent\emph{Form #1}.\quad}
\providecommand{\iftag}[1]{#1}
\providecommand{\tagitem}[1]{\item[#1]}
\providecommand{\Elim}[1]{}
\providecommand{\Intro}[1]{}
\providecommand{\LeftR}[1]{}
\providecommand{\RightR}[1]{}
\providecommand{\TAss}[1]{}
\providecommand{\TRule}[1]{}

\makeatother
\providecommand{\olpath}{.}

% ===================== OLP compatibility layer (2) =====================
% --- second compatibility block: notation and shorthands used by the frozen OLP units ---
\providecommand{\CutCS}{\mathrm{Cut}}
\providecommand{\LeftSubproofLabel}{\relax}
\providecommand{\RightSubproofLabel}{\relax}
\providecommand{\maeh}{\relax}
\providecommand{\cheight}{\relax}
\providecommand{\ande}{\mathrm{and}_E}
\providecommand{\andi}{\mathrm{and}_I}
\providecommand{\ore}{\mathrm{or}_E}
\providecommand{\ori}{\mathrm{or}_I}
\providecommand{\typeof}{\mathord{:}}
\providecommand{\pValue}[1]{#1}
\providecommand{\Sent}{\mathrm{Sent}}
\providecommand{\cardplus}{\mathbin{+}}
\providecommand{\cardtimes}{\mathbin{\cdot}}
\providecommand{\cardexpo}{\mathbin{\hat{\ }}}
\providecommand{\cardnless}{\not<}
\providecommand{\cardle}{\leq}
\providecommand{\funfromto}[3]{#1\colon #2 \to #3}
\providecommand{\funimage}[1]{\operatorname{im}(#1)}
\providecommand{\fdefined}{\mathord{\downarrow}}
\providecommand{\canonord}{\operatorname{ord}}
\providecommand{\ordtimes}{\mathbin{\cdot}}
\providecommand{\ordeq}[2]{#1 = #2}
\providecommand{\ordneq}[2]{#1 \neq #2}
\providecommand{\ordtype}[1]{\operatorname{ord}(#1)}
\providecommand{\ordsucc}{\mathord{+1}}
\providecommand{\ordunion}{\cup}
\providecommand{\onesphere}{S^1}
\providecommand{\rotationsgroup}{\mathrm{SO}(2)}
\providecommand{\equivrep}[1]{[#1]}
\providecommand{\rlexless}{\prec}
\providecommand{\ZFC}{\mathrm{ZFC}}
\providecommand{\ZF}{\mathrm{ZF}}
\providecommand{\Rat}{\mathbb{Q}}
\providecommand{\Real}{\mathbb{R}}
\providecommand{\PosInt}{\mathbb{Z}^{+}}
\providecommand{\Id}{\mathrm{id}}
\providecommand{\Bin}{\mathrm{Bin}}
\providecommand{\pto}{\rightharpoonup}
\providecommand{\dom}{\operatorname{dom}}
\providecommand{\ran}{\operatorname{ran}}
\providecommand{\len}{\operatorname{len}}
\providecommand{\concat}{\mathbin{{}^\frown}}
\providecommand{\card}[1]{\lvert #1 \rvert}
\providecommand{\Pow}{\mathcal{P}}
\providecommand{\Setabs}[2]{\{\, #1 \mid #2 \,\}}
\providecommand{\comp}{\circ}
\providecommand{\Char}{\chi}
\providecommand{\Prf}{\operatorname{Prf}}
\providecommand{\Prov}{\operatorname{Prov}}
\providecommand{\Proj}{\operatorname{proj}}
\providecommand{\Gn}{\operatorname{gn}}
\providecommand{\gn}{\operatorname{gn}}
\providecommand{\Value}[1]{\operatorname{val}(#1)}
\providecommand{\Deduce}{\mathrm{D}}
\providecommand{\Dual}{\mathrm{d}}
\providecommand{\Ax}{\mathrm{Ax}}
\providecommand{\pSat}[3]{#1 \models #2\,[#3]}
\providecommand{\pAssign}{\sigma}
\providecommand{\bforall}{\forall}
\providecommand{\bexists}{\exists}
\providecommand{\liff}{\mathbin{\leftrightarrow}}
\providecommand{\FV}{\operatorname{FV}}
\providecommand{\UnaryInf}{\ensuremath{\mid}}
\providecommand{\BinaryInf}{\ensuremath{\mid}}
\providecommand{\depth}{\operatorname{depth}}
\providecommand{\oliflabeldef}[3]{#1}
\providecommand{\TMblank}{\sqcup}
\providecommand{\TMstroke}{\mid}
\providecommand{\TMstay}{\mathrm{S}}
\providecommand{\TMright}{\mathrm{R}}
\providecommand{\TMleft}{\mathrm{L}}
\providecommand{\TMendtape}{\triangleright}
\providecommand{\umin}{\mu}
\providecommand{\Sum}{\sum}
\providecommand{\citep}[1]{\cite{#1}}
\providecommand{\citet}[1]{\cite{#1}}
\providecommand{\citeauthor}[1]{\cite{#1}}
\providecommand{\citeyear}[1]{\cite{#1}}

% ===================== OLP compatibility layer (3) =====================
% --- third compatibility block: remaining notation seen in the frozen OLP units ---
\providecommand{\iso}{\cong}
\providecommand{\nszero}{\mathrm{0}}
\providecommand{\nssucc}{\mathrm{S}}
\providecommand{\nsplus}{\mathbin{+}}
\providecommand{\nssless}{<}
\providecommand{\nsless}{<}
\providecommand{\PL}{\mathrm{PL}}
\providecommand{\Nec}{\mathord{\Box}}
\providecommand{\RK}{\mathrm{RK}}
\providecommand{\ST}{\mathrm{ST}}
\providecommand{\SSubst}[3]{#1[#2/#3]}
\providecommand{\subst}[3]{#1[#2/#3]}
\providecommand{\tagfalse}{\mathrm{F}}
\providecommand{\cutr}{\mathrm{Cut}}
\providecommand{\cardneq}{\neq}
\providecommand{\issucc}[1]{\operatorname{succ}(#1)}
\providecommand{\tup}[1]{\langle #1 \rangle}
\providecommand{\Var}{\mathrm{Var}}
\providecommand{\Bew}{\operatorname{Bew}}
\providecommand{\Con}{\operatorname{Con}}
\providecommand{\Cons}{\operatorname{Cons}}
\providecommand{\pDeduce}{\mathrm{D}}
\providecommand{\psat}{\models}
\providecommand{\den}[1]{\llbracket #1 \rrbracket}
\providecommand{\Sem}[2]{\llbracket #1 \rrbracket_{#2}}
\providecommand{\DomainOf}[1]{\lvert #1 \rvert}
\providecommand{\Form}{\mathrm{Form}}
\providecommand{\Terms}{\mathrm{Term}}
\providecommand{\Atom}{\mathrm{Atom}}
\providecommand{\FVars}[1]{\operatorname{FV}(#1)}
\providecommand{\fv}[1]{\operatorname{fv}(#1)}
\providecommand{\tuplep}[1]{\langle #1 \rangle}
\providecommand{\restrict}[2]{#1 \upharpoonright #2}
\providecommand{\sto}{\rightarrow}
\providecommand{\FmlaSet}[1]{\mathrm{Fmla}_{#1}}
\providecommand{\Thm}{\mathrm{Thm}}
\providecommand{\Eval}{\operatorname{val}}
\providecommand{\blank}{\sqcup}
\providecommand{\blankTape}{\sqcup}
\providecommand{\Consis}{\operatorname{Con}}
\providecommand{\drs}{\mathrm{d}}
\providecommand{\VRule}{\mid}
\providecommand{\Bet}{\mathrm{B}}
\providecommand{\Seg}[1]{\mathrm{Seg}_{#1}}
\providecommand{\Enums}{\mathrm{En}}
\providecommand{\Z}{\mathbb{Z}}
\providecommand{\N}{\mathbb{N}}
\providecommand{\abs}[1]{\lvert #1 \rvert}
\providecommand{\floor}[1]{\lfloor #1 \rfloor}
\providecommand{\ceil}[1]{\lceil #1 \rceil}

% ===================== OLP compatibility layer (4) =====================
% --- fourth compatibility block: harvested from the per-candidate compile logs ---
\providecommand{\beredpar}{\mathrel{\rightarrow_\beta}}
\makeatletter
\def\bered{\@ifnextchar[{\bered@i}{\mathrel{\rightarrow_\beta}}}
\def\bered@i[#1]{\mathrel{\rightarrow_\beta}}
\def\bred{\@ifnextchar[{\bred@i}{\mathrel{\rightarrow_\beta}}}
\def\bred@i[#1]{\mathrel{\rightarrow_\beta}}
\makeatother
\providecommand{\beredone}{\mathrel{=_\beta}}
\providecommand{\bredone}{\mathrel{=_\beta}}
\providecommand{\xredone}{\mathrel{=_\beta}}
\providecommand{\aconvone}{\mathrel{=_\alpha}}
\providecommand{\cardless}{<}
\providecommand{\tagrefs}[1]{#1}
\providecommand{\ordplus}{\mathbin{+}}
\providecommand{\funrestrictionto}[2]{#1 \upharpoonright #2}
\providecommand{\fundefined}{\mathord{\uparrow}}
\providecommand{\Trm}{\mathrm{Trm}}
\providecommand{\OProv}{\operatorname{Prov}}
\providecommand{\OPrf}{\operatorname{Prf}}
\providecommand{\aeq}{\mathrel{=}}
\providecommand{\mClass}[1]{[#1]}
\providecommand{\equivclass}[1]{[#1]}
\providecommand{\pheight}{\operatorname{ht}}
\providecommand{\phy}{\operatorname{ht}}
\providecommand{\Cut}{\mathrm{Cut}}
\providecommand{\QuantRank}{\operatorname{qr}}
\providecommand{\elemequiv}{\mathrel{\equiv}}
\providecommand{\PIso}{\mathrm{PIso}}
\providecommand{\Theory}{\mathrm{Th}}
\providecommand{\supstrict}{\subsetneq}
\providecommand{\trcl}{\operatorname{trcl}}
\providecommand{\cref}[1]{\ref{#1}}
\providecommand{\Discharge}{\mathrm{Dis}}
\providecommand{\Weakening}{\mathrm{Wk}}
\providecommand{\ZFminus}{\mathrm{ZF}^{-}}
\makeatletter
\renewcommand{\Dual}{\@ifnextchar[{\Dual@i}{\ifmmode\mathrm{d}\else$\mathrm{d}$\fi}}
\def\Dual@i[#1]{\ifmmode\mathrm{d}[#1]\else$\mathrm{d}[#1]$\fi}
\makeatother
\providecommand{\RightR}{\mathrm{R}}
\providecommand{\LeftR}{\mathrm{L}}
\providecommand{\Intro}{\mathrm{I}}
\providecommand{\Axiom}{\mathrm{Ax}}
\providecommand{\formula}[1]{#1}
\providecommand{\stared}{\ast}
\providecommand{\oset}[2]{\{\, #1 \mid #2 \,\}}
\providecommand{\setof}[2]{\{\, #1 \mid #2 \,\}}
\providecommand{\pow}[1]{\mathcal{P}(#1)}
\providecommand{\ordpair}[2]{\langle #1, #2 \rangle}
\providecommand{\ordtriple}[3]{\langle #1, #2, #3 \rangle}
\providecommand{\union}{\cup}
\providecommand{\inter}{\cap}
\providecommand{\seq}[1]{\langle #1 \rangle}
\providecommand{\SuchThat}{\mid}
\providecommand{\where}{\mid}
\providecommand{\TVar}{\mathrm{TVar}}
\providecommand{\Bool}{\mathrm{Bool}}
\providecommand{\Val}{\operatorname{Val}}
\providecommand{\TL}[1]{\mathrm{TL}}
\providecommand{\xstared}{\ast}

% ===================== OLP compatibility layer (5) =====================
% --- fifth compatibility block: rescues the quarantined OLP candidates ---
\providecommand{\ident}{\operatorname{id}}
\providecommand{\setrank}{\operatorname{rank}}
\providecommand{\cardfont}{\mathfrak{c}}
\providecommand{\Complement}{\overline}
\providecommand{\QR}{\mathrm{QR}}
\providecommand{\becd}{\ensuremath{\ulcorner}}
\providecommand{\disjointsum}{\mathbin{\oplus}}
\providecommand{\cardeq}[2]{#1 = #2}
\providecommand{\xred}{\mathrel{\rightarrow}}
% --- sixth compatibility block: names reported by real compile logs ---
\providecommand{\ORefut}{\operatorname{Refut}}
\providecommand{\ORProv}{\operatorname{Prov}}
\providecommand{\cutrank}{\operatorname{cr}}
\providecommand{\Refut}{\operatorname{Refut}}
\providecommand{\DischargeRule}{\mathrm{Dis}}
\providecommand{\Contraction}{\mathrm{C}}
\providecommand{\bottomAlignProof}{\relax}
\providecommand{\FalseCl}{\bot}
\providecommand{\red}{\mathrel{\rightarrow}}
\providecommand{\varAssign}{\sigma}
\providecommand{\defiff}{\mathrel{\leftrightarrow}}
\providecommand{\FalseInt}{\bot}
\providecommand{\fn}{\operatorname{fn}}
\providecommand{\dotted}{\cdot}

% --- tag-macro semantic fix (2026-10-13, final): \iftag CHOOSES a branch and the OLP proof body is
% always the SECOND group (\iftag{tag}{body}, {tag}{body}{c}, {tag}{body}{c}{d}); a definition that
% prints the tag literal drops the proof body. Bodies contain blank lines, hence \long.
\makeatletter
\providecommand{\iftag}{}
\long\def\iftag#1#2{\@ifnextchar\bgroup{\iftag@three{#2}}{\iftag@stop{#2}}}
\long\def\iftag@three#1#2{\iftag@stop{#1}\@ifnextchar\bgroup{\iftag@eat}{\relax}}
\long\def\iftag@eat#1{}
\long\def\iftag@stop#1{#1}
\long\def\oliflabeldef#1#2{\@ifnextchar\bgroup{\@gobble}{\iftag@stop{#2}}}
\providecommand{\tagrefs}[1]{}
\providecommand{\settag}[1]{}
\makeatother
\begin{document}

% ==== OLP file: content/lambda-calculus/lambda-definability/arithmetical-functions.tex ====
% Part: lambda-calculus
% Chapter: lambda-definablity
% Section: arithmetical-functions



\olfileid{lam}{rep}{arf}
\olsection{\usetoken{S}{lambda definable} Arithmetical Functions}

\begin{prop}
  \ollabel{prop:succ-ld}
  The successor function~$\Succ$ is !!{lambda definable}.
\end{prop}

\begin{proof}
A term that !!{lambda define}s the successor function is
\begin{align*}
  \fn{Succ} & \ident \lambd[a][\lambd[fx][f ({a} f x)]].
  \intertext{Given our conventions, this is short for}
  \fn{Succ} & \ident \lambd[a][\lambd[f][\lambd[x][(f (({a} f) x))]]].
  \intertext{$\fn{Succ}$ is a function that accepts as argument a
    number~$a$, and evaluates to another function, $\lambd[fx][f ({a}
      f x)]$. That function is not itself a Church numeral. However,
    if the argument $a$ is a Church numeral, it reduces to one. Consider:}
   (\lambd[a][\lambd[fx][f ({a} f x)]])\,\num{n} & \redone
   \lambd[fx][f ({\num{n}} f x)].
   \intertext{The embedded term $\num{n}fx$ is a redex, since
     $\num{n}$ is $\lambd[fx][f^nx]$. So $\num{n}fx \redone f^nx$ and
     so, for the entire term we have}
   \fn{Succ}\, \num n & \red \lambd[fx][f(f^n(x))],
\end{align*}
i.e., $\num{n+1}$.
\end{proof}

\begin{ex}
  Let's look at what happens when we apply $\fn{Succ}$ to~$\num{0}$,
  i.e., $\lambd[fx][x]$. We'll spell the terms out in full:
  \begin{align*}
    \fn{Succ}\, \num{0} & \ident (\lambd[a][\lambd[f][\lambd[x][(f (({a} f) x))]]])(\lambd[f][\lambd[x][x]])\\
    & \redone \lambd[f][\lambd[x][(f (({(\lambd[f][\lambd[x][x]])} f) x))]] \\
    & \redone \lambd[f][\lambd[x][(f ({(\lambd[x][x])} x))]] \\
    & \redone \lambd[f][\lambd[x][(f x)]] \ident \num{1}\\
  \end{align*}
\end{ex}

\begin{prob}
  The term 
  \begin{align*}
    \fn{Succ}' & \ident \lambd[{n}][\lambd[fx][{n} f (f x)]]
  \end{align*}
  !!{lambda define}s the successor function. Explain why.
\end{prob}

\begin{prop}
  \ollabel{prop:add-ld}
  The addition function~$\Add$ is !!{lambda definable}.
\end{prop}

\begin{proof}
Addition is !!{lambda define}d by the terms  
\begin{align*}
  \fn{Add} & \ident \lambd[{a}{b}][\lambd[fx][{a} f ({b} f x)]]
\intertext{or, alternatively,}
  \fn{Add}' & \ident \lambd[{a}{b}][{a}\, \fn{Succ} \, {b}].
\intertext{The first addition works as follows: $\fn{Add}$ first
  accept two numbers ${a}$ and ${b}$. The result is a function that
  accepts $f$ and $x$ and returns $af(bfx)$. If $a$ and $b$ are Church
  numerals $\num{n}$ and $\num{m}$, this reduces to $f^{n+m}(x)$,
  which is identical to $f^{n}(f^{m}(x))$. Or, slowly:}
  (\lambd[{a}{b}][\lambd[fx][{a} f ({b} f x)]])\num n\,\num m & \redone
  \lambd[fx][\num{n}\, f (\num {m}\, f x)] \\
  & \redone \lambd[fx][\num{n}\, f (f^m x)] \\
  & \redone \lambd[fx][f^n (f^m x)] \ident \num {n+m}.
\intertext{The second representation of addition $\fn{Add'}$ works
  differently: Applied to two Church numerals $\num{n}$ and
  $\num{m}$,}
\fn{Add}' \num n \,\num m
& \redone \num{n}\, \fn{Succ}\, \num{m}.
\intertext{But $\num{n} f x$ always reduces to $f^n(x)$. So,}
  \num{n}\,  \fn{Succ}\, \num{m} & \red \fn{Succ}^n(\num{m}).
\end{align*}
And since $\fn{Succ}$ !!{lambda define}s the successor function, and the
successor function applied $n$ times to~$m$ gives $n+m$, this in turn
reduces to~$\num{n+m}$.
\end{proof}

\begin{prop}
  \ollabel{prop:mult-ld}
  Multiplication is !!{lambda definable} by the term
  \[
  \fn{Mult} \ident \lambd[ab][\lambd[fx][a (b f) x]]
  \]
\end{prop}

\begin{proof}
  To see how this works, suppose we apply $\fn{Mult}$ to Church numerals
  $\num{n}$ and $\num{m}$: $\fn{Mult} \, \num{n} \, \num{m}$ reduces to
  $\lambd[fx][\num{n}(\num{m}\, f)x]$.  The term $\num{m} f$ defines a
  function which applies $f$ to its argument $m$ times. Consequently,
  $\num{n} (\num{m} f) x$ applies the function ``apply $f$ $m$ times''
  itself $n$ times to~$x$. In other words, we apply $f$ to $x$, $n\cdot
  m$ times. But the resulting normal term is just the Church numeral
  $\num{nm}$.
\end{proof}

\begin{editorial}
We can actually simplify this term further by $\eta$-reduction:
\[
  \fn{Mult} \ident \lambd[ab][\lambd[f][a (b f)]].
\]

But then we first have to explain $\eta$-reduction.
\end{editorial}

\begin{prob}
Multiplication can be !!{lambda define}d by the term
\[
  \fn{Mult}' \ident \lambd[ab][a (\fn{Add}\, a) \num{0}].
\]
Explain why this works.
\end{prob}

The definition of exponentiation as a $\lambd$-term is
surprisingly simple:
\[
  \fn{Exp} \ident \lambd[be][e b].
\]
The first argument $b$ is the base and the second~$e$ is the exponent.
Intuitively, $e f$ is $f^e$ by our encoding of numbers. If you find it
hard to understand, we can still define exponentiation also by
iterated multiplication:
\[
  \fn{Exp}' \ident \lambd[be][e (\fn{Mult}\, b) \num{1}].
\]

Predecessor and subtraction on Church numeral is not as simple as we
might think: it requires encoding of pairs.


% ==== OLP file: content/lambda-calculus/lambda-definability/fixpoints.tex ====
% Part: lambda-calculus
% Chapter: lambda-definability
% Section: fixpoints



\olfileid{lam}{ldf}{fp}

\olsection{Fixpoints}

Suppose we wanted to define the factorial function by recursion as a
term $\fn{Fac}$ with the following property:
\[
\fn{Fac} \ident \lambd[n][\fn{IsZero}\, n\, \num 1 (\fn{Mult}\, n (\fn{Fac}(\fn{Pred} \, n)))]
\]
That is, the factorial of $n$ is $1$ if $n = 0$, and $n$ times the
factorial of $n-1$ otherwise.  Of course, we cannot define the term
$\fn{Fac}$ this way since $\fn{Fac}$ itself occurs in the right-hand
side. Such recursive definitions involving self-reference
are not part of the lambda calculus. Defining a term, e.g., by
\[
\fn{Mult} \ident \lambd[ab][a (\fn{Add}\, a) 0]
\]
only involves previously defined terms in the right-hand side, such as
$\fn{Add}$. We can always remove $\fn{Add}$ by replacing it with its
defining term.  This would give the term $\fn{Mult}$ as a pure lambda
term; if $\fn{Add}$ itself involved defined terms (as, e.g.,
$\fn{Add}'$ does), we could continue this process and finally arrive at
a pure lambda term.

However this is not true in the case of recursive definitions like the
one of $\fn{Fac}$ above. If we replace the occurrence of $\fn{Fac}$ on
the right-hand side with the definition of $\fn{Fac}$ itself, we get:
\begin{align*}
  \fn{Fac} & \ident \lambd[n][\fn{IsZero}\, n\, \num{1}] \\
    & \qquad (\fn{Mult}\, n
  ((\lambd[n][\fn{IsZero} \, n \, \num 1\, (\fn{Mult}\, n\, (\fn{Fac}
    (\fn{Pred}\, n)))]) (\fn{Pred}\, n)))
\end{align*}
and we still haven't gotten rid of $\fn{Fac}$ on the right-hand
side. Clearly, if we repeat this process, the definition keeps growing
longer and the process never results in a pure lambda term. Thus this
way of defining factorial (or more generally recursive functions)
is not feasible.

The recursive definition does tell us something, though: If $f$ were a
term representing the factorial function, then the term
\[
\fn{Fac}' \ident \lambd[g][\lambd[n][\fn{IsZero} \, n \, \num 1\, (\fn{Mult} \, n\, (g (\fn{Pred} n)))]]
\]
applied to the term $f$, i.e., $\fn{Fac}'\,f$, also represents the
factorial function.  That is, if we regard $\fn{Fac}'$ as a function
accepting a function and returning a function, the value of
$\fn{Fac'}\, f$ is just~$f$, provided $f$ is the factorial. A
function~$f$ with the property that $\fn{Fac}' \, f \equal[\beta] f$ is called
a \emph{fixpoint} of $\fn{Fac}'$. So, the factorial is a fixpoint of
$\fn{Fac}'$.

There are terms in the lambda calculus that compute the fixpoints of a
given term, and these terms can then be used to turn a term like
$\fn{Fac}'$ into the definition of the factorial. 

\begin{defn}
  \ollabel{defn:Turing-Y}
  The \emph{Y-combinator} is the term:
  \[
  Y \ident (\lambd[ux][x(uux)])(\lambd[ux][x(uux)]).
  \]
\end{defn}

\begin{thm}
  $Y$ has the property that $Yg \red g(Yg)$ for any term $g$. Thus,
  $Yg$ is always a fixpoint of~$g$.
\end{thm}

\begin{proof}
  Let's abbreviate $(\lambd[ux][x(uux)])$ by~$U$, so that $Y \ident UU$. Then
  \begin{align*}
    Y g &\ident (\lambd[ux][x(uux)])U\,g \\
    &\red (\lambd[x][x(UUx)])g\\
    & \red g(UUg) \ident g(Yg).
  \end{align*}
  Since $g(Yg)$ and $Yg$ both reduce to $g(Yg)$, $g(Yg) \equal[\beta]
  Yg$, so $Yg$ is a fixpoint of~$g$.
\end{proof}

Of course, since $Yg$ is a redex, the reduction can continue indefinitely:
\begin{align*}
  Y g &\red g (Y g) \\
    &\red g (g (Y g)) \\
    &\red g(g (g (Y g)))\\
    &\ldots
\end{align*}
So we can think of $Yg$ as $g$ applied to itself infinitely many
times. If we apply~$g$ to it one additional time, we---so to
speak---aren't doing anything extra; $g$~applied to~$g$ applied
infinitely many times to~$Yg$ is still $g$~applied to~$Yg$ infinitely
many times.

Note that the above sequence of $\beta$-reduction steps starting
with~$Yg$ is infinite. So if we apply $Yg$ to some term, i.e.,
consider $(Yg)N$, that term will also reduce to infinitely many
different terms, namely $(g(Yg))N$, $(g(g(Yg)))N$, \dots. It is
nevertheless possible that some \emph{other} sequence of reduction
steps does terminate in a normal form.

Take the factorial for instance. Define $\fn{Fac}$ as $Y\, \fn{Fac}'$
(i.e., a fixpoint of $\fn{Fac'}$). Then:
\begin{align*}
  \fn{Fac}\,\num 3
  &\red Y\, \fn{Fac}' \, \num 3 \\
  &\red \fn{Fac}' (Y \,\fn{Fac}') \, \num 3 \\
  &\ident (\lambd[x][\lambd[n][\fn{IsZero} \, n\, \num 1\,
      (\fn{Mult}\, n\, (x (\fn{Pred}\, n)))]]) \, \fn{Fac} \, \num 3\\
  & \red \fn{IsZero} \, \num 3\, \num 1\,
      (\fn{Mult}\, \num 3\, (\fn{Fac} (\fn{Pred}\, \num 3))) \\
  &\red  \fn{Mult}\, \num 3 \, (\fn{Fac} \, \num 2).
  \intertext{Similarly,}
  \fn{Fac}\,\num 2
  &\red  \fn{Mult}\, \num 2 \, (\fn{Fac} \, \num 1) \\
  \fn{Fac}\,\num 1
  &\red  \fn{Mult}\, \num 1 \, (\fn{Fac} \, \num 0)
  \intertext{but}
  \fn{Fac}\,\num 0
  &\red \fn{Fac}' (Y \,\fn{Fac}') \, \num 0 \\
  &\ident (\lambd[x][\lambd[n][\fn{IsZero} \, n\, \num 1\,
      (\fn{Mult}\, n\, (x (\fn{Pred}\, n)))]]) \, \fn{Fac} \, \num 0\\
  & \red \fn{IsZero} \, \num 0\, \num 1\,
      (\fn{Mult}\, \num 0\, (\fn{Fac} (\fn{Pred}\, \num 0))).\\
  &\red \num 1.
  \intertext{So together}
  \fn{Fac}\,\num 3
  &\red  \fn{Mult}\, \num 3 \,
  (\fn{Mult} \, \num 2\, (\fn{Mult} \, \num 1\, \num 1)).
\end{align*}

What goes for $\fn{Fac'}$ goes for any recursive definition. Suppose we
have a recursive equation
\begin{align*}
g\,x_1\dots x_n & \equal[\beta] N
\intertext{where $N$ may contain~$g$ and $x_1$, \dots,~$x_n$. Then there is
always a term~$G \ident (Y \lambd[g][\lambd[x_1\dots x_n][N]])$ such
that}
G\,x_1 \dots x_n & \equal[\beta] \Subst{N}{G}{g}.
\intertext{For by the fixpoint theorem,}
G \ident (Y \lambd[g][\lambd[x_1\dots x_n][N]]) & \red \lambd[g][\lambd[x_1\dots x_n][N]](Y \lambd[g][\lambd[x_1\dots x_n][N]])\\
& \ident (\lambd[g][\lambd[x_1\dots x_n][N]])\,G
\intertext{and consequently}
G\,x_1 \dots x_n 
& \red (\lambd[g][\lambd[x_1\dots x_n][N]])\,G\,x_1\dots x_n\\
& \red (\lambd[x_1\dots x_n][\Subst{N}{G}{g}])\,x_1\dots x_n\\
  & \red \Subst{N}{G}{g}.
\end{align*}

The $Y$ combinator of \olref{defn:Turing-Y} is due to Alan
Turing. Alonzo Church had proposed a different version which we'll
call~$Y_C$:
\[
Y_C \ident \lambd[g][(\lambd[x][g(xx)])(\lambd[x][g(xx)])].
\]
Church's combinator is a bit weaker than Turing's in that $Yg
\equal[\beta] g(Yg)$ but not $Yg \bred g(Yg)$. Let $V$ be the term
$\lambd[x][g(xx)]$, so that $Y_C \ident \lambd[g][VV]$. Then
\begin{align*}
  VV & \ident (\lambd[x][g(xx)])V \red g(VV)
  \text{ and thus}\\
  Y_C g & \ident (\lambd[g][VV])g \red VV \red g(VV), \text{ but also}\\
  g(Y_C g) & \ident g((\lambd[g][VV])g) \red g(VV).
\end{align*}
In other words, $Y_Cg$ and $g(Y_Cg)$ reduce to a common term $g(VV)$; so $Y_Cg
\equal[\beta] g(Y_Cg)$. This is often enough for applications.



% ==== OLP file: content/lambda-calculus/lambda-definability/introduction.tex ====
% Part: lambda-calculus
% Chapter: lambda-definability
% Section: introduction



\olfileid{lam}{ldf}{int}
\olsection{Introduction}

At first glance, the lambda calculus is just a very abstract calculus
of expressions that represent functions and applications of them to
others. Nothing in the syntax of the lambda calculus suggests that
these are functions of particular kinds of objects, in particular, the
syntax includes no mention of natural numbers. Its basic
operations---application and lambda abstractions---are operations that
apply to any function, not just functions on natural
numbers.

Nevertheless, with some ingenuity, it is possible to define
arithmetical functions, i.e., functions on the natural numbers, in the
lambda calculus. To do this, we define, for each natural number~$n \in
\Nat$, a special $\lambd$-term~$\num n$, the \emph{Church numeral}
for~$n$. (Church numerals are named for Alonzo Church.)

\begin{defn}
  If $n \in \Nat$, the corresponding \emph{Church numeral} $\num{n}$
  represents~$n$:
  \[
    \num{n} \ident \lambd[fx][f^n(x)]
  \]
  Here, $f^n(x)$ stands for the result of applying $f$ to~$x$ $n$
  times. For example, $\num{0}$ is $\lambd[fx][x]$, and $\num{3}$ is
  $\lambd[fx][f(f(f\,x))]$.
\end{defn}
  
The Church numeral $\num n$ is encoded as a lambda term which
represents a function accepting two arguments $f$ and $x$, and
returns $f^n(x)$. Church numerals are evidently in normal form.

A representation of natural numbers in the lambda calculus is only
useful, of course, if we can compute with them.  Computing with Church
numerals in the lambda calculus means applying a $\lambd$-term~$F$ to
such a Church numeral, and reducing the combined term~$F\, \num n$ to
a normal form. If it always reduces to a normal form, and the normal
form is always a Church numeral~$\num m$, we can think of the output
of the computation as being the number~$m$. We can then think of~$F$
as defining a function $f\colon \Nat \to \Nat$, namely the function
such that $f(n) = m$ iff $F\, \num n \red \num m$. Because of the
Church--Rosser property, normal forms are unique if they exist. So if
$F\, \num n \red \num m$, there can be no other term in normal form,
in particular no other Church numeral, that $F \, \num n$ reduces to.

Conversely, given a function $f\colon \Nat \to \Nat$,
we can ask if there is a term $F$ that defines~$f$ in this way. In
that case we say that $F$ \emph{!!{lambda define}s}~$f$, and that $f$ is
!!{lambda definable}. We can generalize this to many-place and partial
functions.

\begin{defn}
Suppose $f\colon \Nat^k \to \Nat$. We say that a lambda term~$F$
\emph{!!{lambda define}s} $f$ if for all 
$n_0$, \dots,~$n_{k-1}$,
\[
F \, \num{n_0} \, \num{n_1} \dots \num{n_{k-1}} \red \num{f(n_0, n_1, \dots,
  n_{k-1})}
\]
if $f(n_0, \dots, n_{k-1})$ is defined, and $F \, \num{n_0} \,
\num{n_1} \dots \num{n_{k-1}}$ has no normal form otherwise.
\end{defn}

A very simple example are the constant functions. The term $C_k \ident
\lambd[x][\num{k}]$ !!{lambda define}s the function $c_k\colon \Nat \to
\Nat$ such that $c(n) = k$. For $C_k \, \num n \ident
(\lambd[x][\num{k}])\num n \redone \num{k}$ for any~$n$. The identity
function is !!{lambda defined} by $\lambd[x][x]$. More complex
functions are of course harder to define, and often require a lot of
ingenuity. So it is perhaps surprising that every computable function
is !!{lambda definable}. The converse is also true: if a function is
!!{lambda definable}, it is computable.


% ==== OLP file: content/lambda-calculus/lambda-definability/lambda-definability.tex ====
% Part: lambda-calculus
% Chapter: lambda-definability



\olchapter{lam}{rep}{Lambda Definability}

\begin{editorial}
  This chapter is experimental. It needs more explanation, and the
  material should be structured better into definitions and
  propositions with proofs, and more examples.
\end{editorial}

\olimport{introduction}
\olimport{arithmetical-functions}
\olimport{pairs}
\olimport{truth-values}
%\olimport{lists}
\olimport{primitive-recursive-functions}
\olimport{fixpoints}
\olimport{minimization}
\olimport{partial-recursive-functions}
\olimport{lambda-definable-recursive}

\OLEndChapterHook


% ==== OLP file: content/lambda-calculus/lambda-definability/lambda-definable-recursive.tex ====
% Part: lambda-calculus
% Chapter: lambda-definability
% Section: lambda-definable-recursive



\olfileid{lam}{dfl}{ldr}
\olsection{\usetoken{S}{lambda definable} Functions are Recursive}

Not only are all partial recursive functions !!{lambda definable}, the
converse is true, too. That is, all !!{lambda definable} functions are
partial recursive.

\begin{thm}
  \ollabel{thm:lambda-computable} If a partial function $f$ is
  !!{lambda definable}, it is partial recursive.
\end{thm}

\begin{proof}
  We only sketch the proof. First, we arithmetize $\lambd$-terms,
  i.e., systematially assign G\"odel numbers to $\lambd$-terms, using
  the usual power-of-primes coding of sequences. Then we define a
  partial recursive function $\fn{normalize}(t)$ operating on the
  G\"odel number~$t$ of a lambda term as argument, and which returns
  the G\"odel number of the normal form if it has one, or is undefined
  otherwise.  Then define two partial recursive functions
  $\fn{toChurch}$ and $\fn{fromChurch}$ that maps natural numbers to
  and from the G\"odel numbers of the corresponding Church numeral.

  Using these recursive functions, we can define the function~$f$ as a
  partial recursive function. There is a $\lambd$-term~$F$ that
  !!{lambda define}s~$f$. To compute $f(n_1, \dots, n_k)$, first obtain
  the G\"odel numbers of the corresponding Church numerals using
  $\fn{toChurch}(n_i)$, append these to $\Gn{F}$ to obtain the G\"odel
  number of the term $F \num{n_1}\dots\num{n_k}$. Now use
  $\fn{normalize}$ on this G\"odel number. If $f(n_1, \dots, n_k)$ is
  defined, $F \num{n_1}\dots\num{n_k}$ has a normal form (which must
  be a Church numeral), and otherwise it has no normal form (and so
  \[\fn{normalize}(\Gn{F\num{n_1}\dots\num{n_k}})\] is undefined).
  Finally, use $\fn{fromChurch}$ on the G\"odel number of the
  normalized term.
\end{proof}


% ==== OLP file: content/lambda-calculus/lambda-definability/lists.tex ====
% Part: lambda-calculus
% Chapter: lambda-definability
% Section: lists



\olfileid{lam}{ldf}{lst}
\olsection{Lists}

The list $[M_1, M_2, \ldots, M_n]$ can be defined to be:
\[
  [M_1, M_2, \ldots, M_n] \ident \lambd[sz][s M_1 (s M_2 (\ldots (s M_n z)))]
\]
Informally speaking, the term representing a list is the
``accumulator'' function of the list: a function that accepts two
terms; the first is a function that receives accumulated value along
with a new element and returns the new accumulated value; the second
is the initial accumulated value; it accumulates elements one by one
and finally returns the result.

\begin{digress}
  If you are familiar with functional programming you will find the
  definition similar to the \emph{right fold}: a list is defined as
  the right fold function of the elements it contains.
\end{digress}

We can easily define some useful functions based on this encoding; 
\begin{align*}
  \fn{Sum} & \ident
  \lambd[l][l \, (\lambd[xa][\fn{Add} \, x \, a])\,  \num{0}]\\
  \fn{Len} & \ident
  \lambd[l][l \, (\lambd[xa][\fn{Add} \, a \, \num{1}]) \num{0}]
\end{align*}

$\fn{Sum}$ calculates the sum of a list of church numerals. It works
by doing an accumulation on the list, where the initial value is
$\num{0}$ and for each element, return $\fn{Add} \, x\, a$ as the new
value (where $x$ is the current element and $a$ is the accumulated
value). The result is the sum of all elements.

\begin{prob}
  What does $\fn{Len}$ do? Explain.
\end{prob}

\begin{prob}
  Define a function that reverses a list.
\end{prob}


% ==== OLP file: content/lambda-calculus/lambda-definability/minimization.tex ====
% Part: lambda-calculus
% Chapter: lambda-definability
% Section: minimization



\olfileid{lam}{ldf}{min}
\olsection{Minimization}

The general recursive functions are those that can be obtained from
the basic functions $\Zero$, $\Succ$, $\Proj{n}{i}$ by composition,
primitive recursion, and regular minimization. To show that all
general recursive functions are !!{lambda definable} we have to show
that any function defined by regular minimization from !!a{lambda
  definable} function is itself !!{lambda definable}.

\begin{lem}
  \ollabel{lem:min} If $f(x_1, \dots, x_k, y)$ is regular and
  !!{lambda definable}, then $g$~defined by
  \[
  g(x_1, \dots, x_k) = \umin{y}{f(x_1,\dots,x_k, y) = 0}
  \]
  is also !!{lambda definable}.
\end{lem}

\begin{proof}
  Suppose the lambda term~$F$ $\lambda$-defines the regular
  function $f(\vec x, y)$. To !!{lambda define}~$h$ we use a search
  function and a fixpoint combinator:
  \begin{align*}
    \fn{Search} & \ident \lambd[g][\lambd[f\,\vec{x}\,y][
        \fn{IsZero} (f\, \vec{x}\, y)\, y\, (g\, \vec{x} (\fn{Succ}\, y)]]\\
    H & \ident \lambd[\vec x][(Y \, \fn{Search}) F\, \vec{x}\, \num{0}],
  \end{align*}
  where $Y$ is any fixpoint combinator. Informally speaking,
  $\fn{Search}$ is a self-referencing function: starting with~$y$,
  test whether $f\, \vec x\, y$ is zero: if so, return~$y$, otherwise call
  itself with $\fn{Succ}\, y$. Thus $(Y \, \fn{Search}) F
  \num{n_1}\dots\num{n_k}\,\num{0}$ returns the least~$m$ for which $f(n_1,
  \dots, n_k, m) = 0$.
  
  Specifically, observe that
  \begin{align*}
    (Y \, \fn{Search}) F \num{n_1}
    \dots\num{n_k}\, \num{m} & \red \num{m}
    \intertext{if $f(n_1, \dots,
      n_k, m) = 0$, or}
    & \red (Y \, \fn{Search}) F\, \num{n_1} \dots\num{n_k}\, \num{m+1}
    \intertext{otherwise. Since $f$ is regular, $f(n_1, \dots, n_k, y)
      = 0$ for some $y$, and so}
    (Y \, \fn{Search}) F \num{n_1} \dots\num{n_k}\,\num{0}
    & \red \num{h(n_1, \dots, n_k)}.
    \end{align*}
\end{proof}


\begin{prop}
  Every general recursive function is !!{lambda definable}.
\end{prop}

\begin{proof}
 By \olref[prf]{lem:basic}, all basic functions are
 !!{lambda definable}, and by \olref[prf]{lem:comp},
 \olref[prf]{lem:prim}, and \olref{lem:min}, the !!{lambda definable}
 functions are closed under composition, primitive recursion, and
 regular minimization.
\end{proof}


% ==== OLP file: content/lambda-calculus/lambda-definability/pairs.tex ====
% Part: lambda-calculus
% Chapter: lambda-definability
% Section: pairs



\olfileid{lam}{ldf}{pai}
\olsection{Pairs and Predecessor}

\begin{defn}
The pair of $M$ and $N$ (written $\tuple{M,N}$) is defined as follows:
\[
\tuple{M,N} \ident \lambd[f][fMN].
\]
\end{defn}
  
Intuitively it is a function that accepts a function, and applies that
function to the two elements of the pair. Following this idea we have
this constructor, which takes two terms and returns the pair containing
them:
\[
  \fn{Pair} \ident \lambd[mn][\lambd[f][fmn]]
\]
Given a pair, we also want to recover its elements.
For this we need two access functions, which accept a pair as argument and
return the first or second elements in it:
\begin{align*}
  \fn{Fst} & \ident \lambd[p][p(\lambd[mn][m])]\\
  \fn{Snd} & \ident \lambd[p][p(\lambd[mn][n])]
\end{align*}

\begin{prob}
  Explain why the access functions $\fn{Fst}$ and $\fn{Snd}$ work.
\end{prob}

Now with pairs we can !!{lambda define} the predecessor function:
\[
  \fn{Pred} \ident \lambd[n][\fn{Fst}(n (\lambd[p][\tuple{\fn{Snd}\, {p}, \fn{Succ}(\fn{Snd}\, {p})}]) \tuple{\num 0, \num 0})]
\]
Remember that $\num n\, f x$ reduces to $f^{n}(x)$; in this
case $f$ is a function that accepts a pair $p$ and returns a new
pair containing the second component of $p$ and the successor of the
second component; $x$ is the pair $\tuple{0,0}$. Thus, the
result is $\tuple{0,0}$ for $n=0$, and $\tuple{\num{n-1}, \num n}$
otherwise. $\fn{Pred}$ then returns the first component of the result.

Subtraction can be defined as $\fn{Pred}$ applied to $a$, $b$ times:
\[
\fn{Sub} \ident \lambd[ab][b \fn{Pred}\, a].
\]
    

% ==== OLP file: content/lambda-calculus/lambda-definability/partial-recursive-functions.tex ====
% Part: lambda-calculus
% Chapter: lambda-definability
% Section: partial-recursive-functions



\olfileid{lam}{ldf}{par}
\olsection{Partial Recursive Functions are \usetoken{S}{lambda definable}}

Partial recursive functions are those obtained from the basic
functions by composition, primitive recursion, and unbounded
minimization. They differ from general recursive function in that the
functions used in unbounded search are not required to be regular.
Not requiring regularity means that functions defined by minimization
may sometimes not be defined. 

At first glance it might seem that the same methods used to show that
the (total) general recursive functions are all !!{lambda definable} can
be used to prove that all partial recursive functions are
!!{lambda definable}.  For instance, the composition of $f$ with $g$ is
!!{lambda define}d by $\lambd[x][F (G x)]$ if $f$ and $g$ are
!!{lambda define}d by terms $F$ and~$G$, respectively. However, when the
functions are partial, this is problematic. When $g(x)$ is undefined,
meaning $G x$ has no normal form. In most cases this means that $F (G
x)$ has no normal forms either, which is what we want.  But consider
when $F$ is $\lambd[x][\lambd[y][y]]$, in which case $F (G x)$ does
have a normal form ($\lambd[y][y]$). 

This problem is not insurmountable, and there are ways to
!!{lambda define} all partial recursive functions in such a way that
undefined values are represented by terms without a normal form.
These ways are, however, somewhat more complicated and less intuitive
than the approach we have taken for general recursive functions. We
record the theorem here without proof:

\begin{thm}
  All partial recursive functions are !!{lambda definable}.
\end{thm}



% ==== OLP file: content/lambda-calculus/lambda-definability/primitive-recursive-functions.tex ====
% Part: lambda-calculus
% Chapter: lambda-definability
% Section: primitive-recursive-functions



\olfileid{lam}{ldf}{prf}
\olsection{Primitive Recursive Functions are \usetoken{S}{lambda definable}} 

Recall that the primitive recursive functions are those that can be
defined from the basic functions $\Zero$, $\Succ$, and $\Proj{n}{i}$
by composition and primitive recursion.

\begin{lem}
  \ollabel{lem:basic}
  The basic primitive recursive functions $\Zero$, $\Succ$, and
  projections~$\Proj{n}{i}$ are !!{lambda definable}.
\end{lem}

\begin{proof}
They are !!{lambda define}d by the following terms:
\begin{align*}
  \fn{Zero} & \ident \lambd[a][\lambd[fx][x]]\\
  \fn{Succ} & \ident \lambd[a][\lambd[fx][f (a f x)]] \\
  \fn{Proj}^n_i & \ident \lambd[x_0\dots x_{n-1}][x_i]
\end{align*}
\end{proof}

\begin{lem}
  \ollabel{lem:comp} Suppose the $k$-ary function $f$, and $n$-ary
  functions $g_0, \dots, g_{k-1}$, are !!{lambda definable} by terms
  $F$, $G_0$, \dots, $G_k$, and $h$ is defined from them by composition.
  Then $H$ is !!{lambda definable}.
\end{lem}

\begin{proof}
  $h$ can be !!{lambda define}d by the term
  \[
  H \ident \lambd[x_0 \dots x_{n-1}][F\, (G_0 x_0 \dots
    x_{n-1}) \dots (G_{k-1} x_0 \dots x_{n-1})]
  \]
  We leave verification of this fact as an exercise.
\end{proof}

\begin{prob}
  Complete the proof of \olref[lam][ldf][prf]{lem:comp} by showing
  that $H\num{n_0}\dots\num{n_{n-1}} \red \num{h(n_0, \dots,
    n_{n-1})}$.
\end{prob}

Note that \olref{lem:comp} did not require that $f$ and $g_0$,
\dots,~$g_{k-1}$ are primitive recursive; it is only required that
they are total and !!{lambda definable}.

\begin{lem}\ollabel{lem:prim}
  Suppose $f$ is an $n$-ary function and~$g$ is an $n+2$-ary function,
  they are !!{lambda definable} by terms $F$ and~$G$, and the
  function~$h$ is defined from $f$ and $g$ by primitive
  recursion. Then $h$ is also !!{lambda definable}.
\end{lem}

\begin{proof}
  Recall that $h$ is defined by
  \begin{align*}
    h(x_1, \dots, x_n, 0) &= f(x_1, \dots, x_n)\\
    h(x_1, \dots, x_n, y+1) & = h(x_1, \dots, x_n, y, h(x_1, \dots, x_n, y)).
  \end{align*}
  Informally speaking, the primitive recursive definition iterates the
  application of the function $h$ $y$ times and applies it to $f(x_1,
  \dots, x_n)$. This is reminiscent of the definition of Church
  numerals, which is also defined as a iterator.

  For simplicity, we give the definition and proof for a single
  additional argument~$x$. The function $h$ is !!{lambda define}d by:
  \begin{align*}
    H \ident & \lambd[x][\lambd[y][\fn{Snd} (y
        D \tuple{\num{0}, F x})]]
    \intertext{where }
    D \ident & \lambd[p][\tuple{\fn{Succ} (\fn{Fst}\, p), (G x
          (\fn{Fst}\, p) (\fn{Snd}\, p))}]
  \end{align*}
  The iteration state we maintain is a pair, the first of which is the
  current $y$ and the second is the corresponding value of~$h$. For
  every step of iteration we create a pair of new values of $y$ and
  $h$; after the iteration is done we return the second part of the
  pair and that's the final $h$ value. We now prove this is indeed a
  representation of primitive recursion.

  We want to prove that for any $n$ and $m$, $H\,\num{n}\,\num{m} \red
  \num{h(n,m)}$. To do this we first show that if $D_n \ident
  \Subst{D}{\num{n}}{x}$, then $D_n^m \tuple{\num{0}, F\, \num{n}} \red
  \tuple{\num{m}, \num{h(n, m)}}$ We proceed by induction on~$m$.

  If $m=0$, we want $D_n^0 \tuple{\num{0}, F\, \num{n}} \red
  \tuple{\num{0}, \num{h(n, 0)}}$. But $D_n^0 \tuple{\num{0}, F\,
    \num{n}}$ just is $\tuple{\num{0}, F\, \num{n}}$. Since $F$
  !!{lambda define}s~$f$, this reduces to $\tuple{\num{0},
    \num{f(n)}}$, and since $f(n) = h(n, 0)$, this is $\tuple{\num{0},
    \num{h(n,0)}}$

  Now suppose that $D_n^m \tuple{\num{0}, F\, \num{n}} \red
  \tuple{\num{m}, \num{h(n, m)}}$. We want to
  show that $D_n^{m+1} \tuple{\num{0}, F\, \num{n}} \red
  \tuple{\num{m+1}, \num{h(n, m+1)}}$.
  \begin{align*}
    D_n^{m+1} \tuple{\num{0}, F\, \num{n}}
    & \ident D_n(D_n^{m} \tuple{\num{0}, F\, \num{n}})\\
    & \red D_n\,\tuple{\num{m}, \num{h(n, m)}} \text{\qquad (by IH)}\\
    & \ident (\lambd[p][
      \tuple{
        \fn{Succ} (\fn{Fst}\, p), (G \, \num{n} (\fn{Fst}\, p) (\fn{Snd}\, p))
    }]) \tuple{\num{m}, \num{h(n,m)}}\\
    & \redone
    \tuple{\fn{Succ} (\fn{Fst}\, \tuple{\num{m}, \num{h(n,m)}}),\\
      & \qquad (G \, \num{n} (\fn{Fst}\, \tuple{\num{m}, \num{h(n,m)}}) (\fn{Snd}\, \tuple{\num{m}, \num{h(n,m)}}))}\\
    & \red 
    \tuple{\fn{Succ}\, \num{m}, (G \, \num{n} \,\num{m} \, \num{h(n,m)})}\\
    & \red \tuple{\num{m+1}, \num{g(n, m, h(n, m))}}
  \end{align*}
  Since $g(n, m, h(n, m)) = h(n, m+1)$, we are done.

  Finally, consider
  \begin{align*}
    H\,\num n\,\num m &\ident \lambd[x][\lambd[y][\fn{Snd} (y
        (\lambd[p].\tuple{\fn{Succ} (\fn{Fst}\, p), (G\, x\,
          (\fn{Fst}\, p)\, (\fn{Snd}\, p))}) \tuple{\num{0}, F x})]]\\
    & \qquad\,\num n\, \num m\\
    & \red \fn{Snd} (\num m \,
    \underbrace{(\lambd[p].\tuple{\fn{Succ} (\fn{Fst}\, p), (G \,\num n\,
      (\fn{Fst}\, p) (\fn{Snd}\, p))})}_{D_n} \tuple{\num{0}, F \num n})\\
    & \ident \fn{Snd} (\num m \, D_n\, \tuple{\num{0}, F \num n})\\
    & \red \fn{Snd} \,(D_n^m  \tuple{\num{0}, F \num n})\\
    & \red \fn{Snd} \,\tuple{\num{m}, \num{h(n, m)}} \\
    & \red \num{h(n,m)}. 
  \end{align*}  
\end{proof}


\begin{prop}
  Every primitive recursive function is !!{lambda definable}.
\end{prop}

\begin{proof}
  By \olref{lem:basic}, all basic functions are !!{lambda definable},
  and by \olref{lem:comp} and \olref{lem:prim}, the !!{lambda definable}
  functions are closed under composition and primitive recursion.
\end{proof}



% ==== OLP file: content/lambda-calculus/lambda-definability/truth-values.tex ====
% Part: lambda-calculus
% Chapter: lambda-definability
% Section: truth-values



\olfileid{lam}{ldf}{tvr}
\olsection{Truth Values and Relations}

We can encode truth values in the pure lambda calculus as follows:
\begin{align*}
  \fn{true} & \ident \lambd[x][\lambd[y][x]]\\
  \fn{false} & \ident \lambd[x][\lambd[y][y]]
\end{align*}

Truth values are represented as \emph{selectors}, i.e., functions that
accept two arguments and returning one of them. The truth value
$\fn{true}$ selects its first argument, and $\fn{false}$ its
second. For example, $\fn{true}\, M N$ always reduces to $M$, while
$\fn{false}\, M N$ always reduces to~$N$.

\begin{defn}
We call a relation $R \subseteq \Nat^n$ !!{lambda definable} if there is
a term~$R$ such that
\begin{align*}
  R\, \num{n_1} \dots \num{n_k} & \bred \fn{true}
  \intertext{whenever $R(n_1, \dots, n_k)$ and}
  R\, \num{n_1} \dots \num{n_k} & \bred \fn{false}
\end{align*}
otherwise.
\end{defn}

For instance, the relation $\fn{IsZero} = \{0\}$ which holds of $0$
and $0$ only, is !!{lambda definable} by
\[
  \fn{IsZero} \ident \lambd[n][n (\lambd[x][\fn{false}])\, \fn{true}].
\]
How does it work? Since Church numerals are defined as iterators
(functions which apply their first argument $n$~times to the second),
we set the initial value to be $\fn{true}$, and for every step of
iteration, we return $\fn{false}$ regardless of the result of the last
iteration.  This step will be applied to the initial value $n$
times, and the result will be $\fn{true}$ if and only if the step is
not applied at all, i.e., when $n = 0$.

On the basis of this representation of truth values, we can further
define some truth functions. Here are two, the representations of
negation and conjunction:
\begin{align*}
  \fn{Not} & \ident \lambd[x][x\, \fn{false}\, \fn{true}]\\
  \fn{And} & \ident \lambd[x][\lambd[y][xy \,\fn{false}]]
\end{align*}
The function ``$\fn{Not}$'' accepts one argument, and returns
$\fn{true}$ if the argument is $\fn{false}$, and $\fn{false}$ if the
argument is~$\fn{true}$.  The function ``$\fn{And}$'' accepts two
truth values as arguments, and should return $\fn{true}$ iff both
arguments are~$\fn{true}$. Truth values are represented as selectors
(described above), so when $x$ is a truth value and is applied to two
arguments, the result will be the first argument if $x$ is $\fn{true}$
and the second argument otherwise. Now $\fn{And}$ takes its two
arguments $x$ and $y$, and in return passes $y$ and $\fn{false}$ to
its first argument~$x$. Assuming $x$ is a truth value, the result
will evaluate to~$y$ if $x$ is $\fn{true}$, and to $\fn{false}$ if $x$
is $\fn{false}$, which is just what is desired.

Note that we assume here that only truth values are used as arguments
to $\fn{And}$. If it is passed other terms, the result (i.e., the
normal form, if it exists) may well not be a truth value.

\begin{prob}
  Define the functions $\fn{Or}$ and $\fn{Xor}$ representing the truth
  functions of inclusive and exclusive disjunction using the encoding
  of truth values as $\lambd$-terms.
\end{prob}

\end{document}

\end{document}

\end{document}
